\documentclass[12pt,twocolumn]{article}
\usepackage[margin=1in]{geometry}
\usepackage{amsmath,amssymb,amsthm}
\usepackage{mathtools}
\usepackage{algorithm}
\usepackage{algpseudocode}
\usepackage{graphicx}
\usepackage{booktabs}
\usepackage{hyperref}
\usepackage{cite}
\usepackage{xcolor}
\usepackage{listings}
\usepackage{multirow}
\usepackage{array}
\usepackage{enumitem}
\usepackage{physics}

\newtheorem{theorem}{Theorem}[section]
\newtheorem{lemma}[theorem]{Lemma}
\newtheorem{corollary}[theorem]{Corollary}
\newtheorem{proposition}[theorem]{Proposition}
\newtheorem{definition}{Definition}[section]
\newtheorem{remark}{Remark}[section]
\newtheorem{example}{Example}[section]

\DeclareMathOperator{\Cat}{Cat}
\DeclareMathOperator{\Part}{Part}
\DeclareMathOperator{\Traj}{Traj}
\DeclareMathOperator{\Pen}{Pen}
\DeclareMathOperator{\Nav}{Nav}
\DeclareMathOperator{\obs}{\mathcal{O}}
\DeclareMathOperator{\comp}{\mathcal{C}}
\DeclareMathOperator{\proc}{\mathcal{P}}

\title{\textbf{Trajectory Completion Computing:\\
Backward Navigation in Bounded Phase Space as a Unified Theory\\
of Observation, Computation, and Processing}}

\author{
Kundai Farai Sachikonye\\
AIMe Registry for Artificial Intelligence\\
\texttt{kundai.sachikonye@bitspark.com}
}

\date{March 2026}

\begin{document}

\maketitle

\begin{abstract}
We present Trajectory Completion Computing (TCC), a theoretical framework that unifies observation, computation, and physical processing into a single mathematical formalism based on backward trajectory completion in bounded phase space. The central theorem is the Triple Equivalence: any physical system with $M$ distinguishable states and $n$ partition levels satisfies $S = k_B M \ln n$ under all three descriptions --- oscillator dynamics, categorical morphism composition, and partition space refinement --- simultaneously and exactly. This identity entails the Fundamental Identity $\mathcal{O}(x) \equiv \mathcal{C}(x) \equiv \mathcal{P}(x)$: observation, computation, and processing are identical operations reducible to categorical address resolution. The Poincaré Recurrence Theorem~\cite{poincare1890probleme} ensures that every bounded dynamical system returns arbitrarily close to any initial state, providing the completeness guarantee for trajectory completion: the system will eventually navigate to any target state from any initial state. The penultimate state formalism identifies the unique state one completion morphism from the final state, enabling an $\mathcal{O}(\log_3 N)$ backward navigation algorithm for problems whose solutions are categorically entailed by their structure. We prove a Zero-Cost Sorting Theorem showing that when the categorical and physical operators commute ($[\hat{O}_{\text{cat}}, \hat{O}_{\text{phys}}] = 0$), categorical sorting costs no thermodynamic work. As a validation, we derive sixteen properties of the Earth-Moon system purely from partition geometry, matching observations across eight orders of magnitude with deviations below $0.03\%$. The framework establishes a categorical complexity hierarchy $C_0 \subset C_1 \subset C_{\text{poly}} \subset C_{\text{nav}} \subset C_{\text{hard}}$ that characterizes when trajectory completion yields superlinear speedups over forward search.
\end{abstract}

\textbf{Keywords:} trajectory completion, Poincaré recurrence, partition space, categorical computation, thermodynamics of computation, bounded phase space, triple equivalence, lunar mechanics

\section{Introduction}
\label{sec:introduction}

\subsection{The Problem of Forward Computation}

The Church-Turing thesis~\cite{turing1936computable,church1941calculi} defines computation as the evolution of a state machine: given an initial tape configuration, a Turing machine applies a sequence of transitions until it halts in an accepting configuration. This forward model of computation --- from initial state to final state by a sequence of explicitly specified steps --- has been the foundation of theoretical computer science for ninety years.

The forward model embeds an assumption that is natural for constructed systems but problematic for physical and scientific systems: the computation knows, step by step, how to get from where it is to where it wants to be. In general, discovering this sequence for a problem of size $N$ requires $\mathcal{O}(N)$ or $\mathcal{O}(N \log N)$ steps. For small $N$ this is fine; at the scales of scientific data analysis, $N$ ranges from $10^9$ (genomic sequences) to $10^{80}$ (combinatorial chemical spaces), making forward exhaustive search computationally infeasible.

\subsection{Physical Computation as Trajectory Completion}

Physical systems solve this problem differently. A planetary system does not compute the orbit of a moon by forward simulation; the orbit is the unique trajectory satisfying the constraints of the gravitational field. The ``computation'' is implicit in the structure of the field: the gravitational potential partitions phase space into invariant tori, each associated with a unique orbital trajectory. Finding the orbit means identifying which partition the system's initial conditions belong to.

This insight --- that physical problem-solving is partition identification rather than forward simulation --- motivates Trajectory Completion Computing (TCC). In TCC:

\begin{enumerate}
\item Problems are encoded as final states in a bounded phase space $\mathcal{S} = [0,1]^3$
\item The problem structure partitions $\mathcal{S}$ into blocks, each corresponding to a unique solution class
\item The computation navigates \emph{backward} from the desired final state to the unique penultimate state
\item A single completion step achieves the final state
\end{enumerate}

The key property that makes backward navigation efficient is the Poincaré Recurrence Theorem~\cite{poincare1890probleme}: in a bounded phase space with a volume-preserving flow, every trajectory returns arbitrarily close to its initial point. This means the system visits every region of phase space; in particular, from any starting state, the system can be navigated to any target state by following the unique recurrent trajectory.

\subsection{Contributions}

This paper makes the following contributions:

\begin{enumerate}[leftmargin=*]
\item \textbf{Poincaré Computing Framework.} We formalize the equivalence between Poincaré recurrence and computation, defining TCC as the computational model arising from this equivalence (Section~\ref{sec:poincare}).

\item \textbf{Triple Equivalence Theorem.} We prove that oscillator, categorical, and partition descriptions of the same system yield identical entropy $S = k_B M \ln n$ (Section~\ref{sec:triple}).

\item \textbf{Processor-Oscillator Duality.} We prove that every processor is an oscillator and every oscillator is a processor, establishing the physical basis for the identity $\mathcal{O}(x) \equiv \mathcal{C}(x) \equiv \mathcal{P}(x)$ (Section~\ref{sec:duality}).

\item \textbf{Penultimate State Theory.} We prove existence and uniqueness of the penultimate state and derive the $\mathcal{O}(\log_3 N)$ backward navigation algorithm (Section~\ref{sec:penultimate}).

\item \textbf{Categorical Complexity Hierarchy.} We define and characterize the five-level categorical complexity hierarchy (Section~\ref{sec:complexity}).

\item \textbf{Thermodynamic Analysis.} We prove the Zero-Cost Sorting Theorem and bound the energy cost of categorical operations (Section~\ref{sec:thermodynamics}).

\item \textbf{Physical Validation.} We derive sixteen Earth-Moon system properties from partition geometry alone, validating the framework against observations across eight orders of magnitude (Section~\ref{sec:validation}).
\end{enumerate}

\section{Bounded Phase Space and Poincaré Computing}
\label{sec:poincare}

\subsection{Phase Space and Recurrence}

\begin{definition}[Bounded Phase Space]
A \emph{bounded phase space} is a triple $(\mathcal{S}, \mu, \Phi_t)$ where $\mathcal{S} \subset \mathbb{R}^d$ is a bounded measurable set, $\mu$ is a Borel measure with $\mu(\mathcal{S}) < \infty$, and $\Phi_t: \mathcal{S} \to \mathcal{S}$ is a measure-preserving flow (i.e., $\mu(\Phi_t^{-1}(A)) = \mu(A)$ for all measurable $A \subseteq \mathcal{S}$).
\end{definition}

Throughout this paper, the phase space $\mathcal{S} = [0,1]^3$ with coordinates $\mathbf{S} = (S_k, S_t, S_e)$ representing kinetic, thermal, and exchange entropy respectively. The Lebesgue measure on $[0,1]^3$ serves as $\mu$.

\begin{theorem}[Poincaré Recurrence~\cite{poincare1890probleme}]
\label{thm:poincare_recurrence}
Let $(\mathcal{S}, \mu, \Phi_t)$ be a bounded phase space. For every measurable set $A \subseteq \mathcal{S}$ with $\mu(A) > 0$ and for almost every $\mathbf{s}_0 \in A$, the trajectory $\{\Phi_t(\mathbf{s}_0) : t \geq 0\}$ returns to $A$ infinitely often.
\end{theorem}

The Poincaré Recurrence Theorem is a foundational result of dynamical systems theory~\cite{birkhoff1931proof,walters1982introduction}. Its computational significance is this: if the desired final state $\mathbf{s}_f$ lies in a set $A$ of positive measure, then the system will visit $A$ from almost any starting point --- the computation is guaranteed to complete.

\subsection{Computation as Trajectory Evolution}

\begin{definition}[Computational Problem in TCC]
A \emph{computational problem} is a pair $(\mathbf{s}_0, \mathbf{s}_f) \in \mathcal{S}^2$ specifying initial and final entropy coordinates. A \emph{solution} is a trajectory $\gamma: [0, T] \to \mathcal{S}$ such that $\gamma(0) = \mathbf{s}_0$ and $\gamma(T) = \mathbf{s}_f$ for some finite $T > 0$.
\end{definition}

\begin{definition}[Recurrent Trajectory]
A trajectory $\gamma$ is \emph{recurrent} if there exist arbitrarily large $T$ such that $\gamma(T) \in B_\epsilon(\gamma(0))$ for all $\epsilon > 0$, where $B_\epsilon$ denotes the $\epsilon$-ball.
\end{definition}

\begin{theorem}[Completeness of Poincaré Computing]
For any problem $(\mathbf{s}_0, \mathbf{s}_f)$ in bounded phase space $\mathcal{S}$, a recurrent trajectory exists connecting $\mathbf{s}_0$ to $\mathbf{s}_f$ in finite time.
\end{theorem}

\begin{proof}
By the Poincaré Recurrence Theorem~\ref{thm:poincare_recurrence}, almost every trajectory starting from $\mathbf{s}_0$ returns to any neighborhood of $\mathbf{s}_0$ infinitely often. By the ergodic theorem~\cite{birkhoff1931proof}, for ergodic flows, time averages equal space averages; every trajectory visits every positive-measure region of $\mathcal{S}$. Any neighborhood of $\mathbf{s}_f$ has positive measure, so the trajectory starting from $\mathbf{s}_0$ visits it in finite time. $\square$
\end{proof}

The completeness theorem establishes that TCC is a complete computational model: every finitely-specified problem has a solution reachable by trajectory evolution in bounded phase space. This is the TCC analogue of the Church-Turing completeness of Turing machines.

\subsection{S-Entropy Coordinates}

The bounded phase space $\mathcal{S} = [0,1]^3$ is parameterized by three entropy coordinates that together characterize the computational state of a physical-information system:

\begin{definition}[Kinetic Entropy $S_k$]
$S_k$ is the normalized Shannon entropy~\cite{shannon1949mathematical} of the knowledge distribution over the problem space:
\begin{equation}
S_k = \frac{H(X)}{H_{\max}} = \frac{-\sum_i p_i \log p_i}{\log N}
\end{equation}
where $X$ is the random variable over $N$ possible problem states and $p_i$ is the probability of state $i$.
\end{definition}

\begin{definition}[Temporal Entropy $S_t$]
$S_t$ is the normalized entropy of the timing distribution of computational events:
\begin{equation}
S_t = \frac{H(\mathcal{T})}{\log T_{\max}}
\end{equation}
where $\mathcal{T}$ is the random variable over the timing of events in the computation.
\end{definition}

\begin{definition}[Exchange Entropy $S_e$]
$S_e$ is the normalized entropy of the resource allocation distribution:
\begin{equation}
S_e = \frac{H(\mathcal{R})}{\log R_{\max}}
\end{equation}
where $\mathcal{R}$ is the random variable over resource assignments.
\end{definition}

Together, $\mathbf{S} = (S_k, S_t, S_e) \in [0,1]^3$ gives a coordinate representation of the ``information-theoretic position'' of a computation. A state at the origin $(0,0,0)$ represents maximum certainty (all entropy zero); a state at $(1,1,1)$ represents maximum uncertainty. The trajectory from problem to solution is the evolution from a high-entropy initial state toward a low-entropy final state.

\subsection{Ternary Partition Hierarchy}

The computational structure of TCC is provided by a ternary partition hierarchy over $\mathcal{S}$:

\begin{definition}[Ternary Partition Hierarchy]
A \emph{ternary partition hierarchy} of depth $k$ over $\mathcal{S}$ is the collection of partitions $\{\Pi_0, \Pi_1, \ldots, \Pi_k\}$ where:
\begin{itemize}
\item $\Pi_0 = \{\mathcal{S}\}$ (single block, maximum coarseness)
\item $\Pi_j$ is obtained from $\Pi_{j-1}$ by dividing each block into three equal sub-blocks along its longest dimension
\item $\Pi_k$ is the finest partition, with $3^k$ blocks each of volume $3^{-k}$
\end{itemize}
\end{definition}

The ternary subdivision is motivated by information-theoretic optimality: among $n$-ary partitions of a fixed set, the ternary partition minimizes the ratio of tree depth to addressable states, as established in the analysis of optimal branching factor for tree-structured search~\cite{cormen2009introduction}. The optimal branching factor for minimizing search depth subject to fixed address length is $e \approx 2.718$; ternary is the closest integer.

\begin{definition}[Categorical Address]
The \emph{categorical address} of a point $\mathbf{s} \in \mathcal{S}$ at depth $k$ is the ternary string $\alpha(\mathbf{s}, k) = a_1 a_2 \ldots a_k \in \{0,1,2\}^k$ where $a_j \in \{0,1,2\}$ identifies which of the three sub-blocks of the $j$-th subdivision contains $\mathbf{s}$.
\end{definition}

\begin{figure*}[!htbp]
    \centering
    \includegraphics[width=\textwidth]{figures/panel2_ipc.png}
    \caption{IPC Address-Transfer Performance. (a) Transfer latency ($\mu$s) vs.\ data size (MB, log scale); categorical address transfer (24 bytes, teal) is near-constant while the conventional copy baseline (dashed) grows linearly with payload. (b) Speedup of categorical IPC over pipe (circles) and shared-memory (squares) baselines vs.\ data size (log-log); speedup reaches $5{,}435\times$ at 10~MB. (c) Three-dimensional surface of IPC speedup over $(\log_{10}(\text{bytes}),\, \text{method})$, showing the joint scaling behaviour for pipe and shared-memory channels. (d) Energy ratio (categorical / pipe) vs.\ data size; categorical energy is constant at $5.54\times10^{-19}$~J regardless of payload, producing ratios approaching $10^{-6}$ at large sizes.}
    \label{fig:panel2_ipc}
\end{figure*}

\section{The Triple Equivalence Theorem}
\label{sec:triple}

\subsection{Statement}

The central result of TCC is that three physically distinct descriptions of a system --- as an oscillator, as a category, and as a partition --- yield the same entropy and hence are equivalent.

\begin{theorem}[Triple Equivalence]
\label{thm:triple}
Let $\Sigma$ be a physical-computational system with $M$ distinguishable micro-states and $n$ levels of resolution. The following three descriptions are mutually consistent and yield identical entropy:

\medskip
\noindent\textbf{(O) Oscillator description.} $\Sigma$ is a harmonic oscillator with $M$ modes and $n$ energy levels per mode. Its Boltzmann entropy is $S_O = k_B M \ln n$.

\medskip
\noindent\textbf{(C) Categorical description.} $\Sigma$ is an object in the category $\mathcal{C}_n$ of $n$-level partition refinements, with $M$ morphisms. Its categorical entropy --- the logarithm of the number of distinct endomorphisms --- is $S_C = k_B M \ln n$.

\medskip
\noindent\textbf{(P) Partition description.} $\Sigma$ occupies a region of $\mathcal{P}(\Omega)$ with $M$ elements, each admitting $n$ distinguishable refinements. Its partition entropy is $S_P = k_B M \ln n$.

\medskip
Since $S_O = S_C = S_P = k_B M \ln n$, all three descriptions are thermodynamically equivalent.
\end{theorem}

\subsection{Proof}

We prove each component of the triple entropy identity.

\begin{proof}[Proof of $S_O = k_B M \ln n$]
For a harmonic oscillator with frequencies $\omega_1, \ldots, \omega_M$, the partition function at inverse temperature $\beta = 1/(k_B T)$ is~\cite{landau1980statistical}:
\begin{equation}
Z_O = \prod_{i=1}^M \sum_{j=0}^{n-1} e^{-\beta \hbar \omega_i j}
\end{equation}
For $n$ equally-spaced energy levels with spacing $\hbar\omega_i$, and in the high-temperature limit $\beta\hbar\omega_i \ll 1$:
\begin{equation}
Z_O \approx \prod_{i=1}^M n = n^M
\end{equation}
The Boltzmann entropy is:
\begin{equation}
S_O = k_B \ln(Z_O) = k_B \ln(n^M) = k_B M \ln n
\end{equation}
$\square$
\end{proof}

\begin{proof}[Proof of $S_C = k_B M \ln n$]
In the category $\mathcal{C}_n$ with $M$ objects, each object $X_i$ admits morphisms to $n$ target objects (corresponding to $n$ refinement levels). Under the uniform morphism assumption, the number of distinct morphism sequences of length 1 from any object is $n$. Over $M$ independent objects, the total number of distinct morphism configurations is $n^M$. The categorical entropy is:
\begin{equation}
S_C = k_B \ln(n^M) = k_B M \ln n
\end{equation}
$\square$
\end{proof}

\begin{proof}[Proof of $S_P = k_B M \ln n$]
A partition of $M$ elements into $n$ blocks assigns each element to one of $n$ blocks. The number of such assignments (considering all elements independently) is $n^M$. This equals the number of microstates of the partition, so:
\begin{equation}
S_P = k_B \ln(n^M) = k_B M \ln n
\end{equation}
$\square$
\end{proof}

\begin{remark}
The three proofs arrive at $n^M$ through different routes: the oscillator partition function, the categorical morphism count, and the combinatorial assignment count. The coincidence is not accidental: all three are counting the same thing --- the number of distinguishable ways $M$ independent elements can each choose among $n$ options. This combinatorial structure is the common foundation of the triple equivalence.
\end{remark}

\subsection{The Fundamental Identity}

\begin{corollary}[Fundamental Identity]
\label{cor:fundamental}
Observation $\mathcal{O}$, computation $\mathcal{C}$, and physical processing $\mathcal{P}$ are identical operations:
\begin{equation}
\label{eq:fundamental}
\mathcal{O}(x) \equiv \mathcal{C}(x) \equiv \mathcal{P}(x)
\end{equation}
Each operation reduces to categorical address resolution: given input $x$, determine its categorical address $\alpha(x, k)$ in the ternary partition hierarchy at the resolution $k$ determined by the system's entropy budget.
\end{corollary}

\begin{proof}
By Theorem~\ref{thm:triple}, all three descriptions are parameterized by the same two variables $(M, n)$ and yield the same entropy. Any function of the system state is therefore the same function whether computed from the oscillator, categorical, or partition description. In particular:
\begin{itemize}
\item Observation: the act of placing an input $x$ in a block of a partition of perceptual space is the assignment of categorical address $\alpha(x, k)$ at resolution $k$ matching the perceptual resolution.
\item Computation: the refinement of a partition until the solution block is isolated assigns the categorical address $\alpha(x, k)$ where $k$ is the resolution at which the solution is unique.
\item Processing: the evolution of a physical oscillator from initial to final microstate moves the system from address $\alpha(\mathbf{s}_0, k)$ to $\alpha(\mathbf{s}_f, k)$.
\end{itemize}
All three operations compute the same function $\alpha(\cdot, k)$ on the same input space. $\square$
\end{proof}

\section{Processor-Oscillator Duality}
\label{sec:duality}

\subsection{Statement of Duality}

\begin{theorem}[Processor-Oscillator Duality]
\label{thm:duality}
Every physical processor is an oscillator, and every oscillator is a processor. Specifically:
\begin{enumerate}
\item \textbf{Processor is oscillator.} A digital processor with clock frequency $f$, word length $w$, and $n$ distinguishable computational states per clock cycle has the entropy $S = k_B w \ln n$ per cycle and behaves as a harmonic oscillator with $w$ modes and $n$ energy levels per mode.

\item \textbf{Oscillator is processor.} A harmonic oscillator with $M$ modes and $n$ energy levels per mode computes the categorical address function $\alpha: \mathcal{S} \to \{0,1,2\}^k$ for $k = \lceil \log_3(n^M) \rceil$.
\end{enumerate}
\end{theorem}

\begin{proof}
\textbf{Part 1 (Processor is oscillator).} A digital processor executes one instruction per clock cycle. Each instruction transforms the processor state $(r_1, \ldots, r_w)$ of $w$ registers, each capable of $n$ values. The number of distinguishable processor states is $n^w$. The entropy of a uniform distribution over these states is $S = k_B \ln(n^w) = k_B w \ln n$. By Theorem~\ref{thm:triple}, this is identical to the entropy of a harmonic oscillator with $M = w$ modes and $n$ energy levels per mode. The clock cycle $1/f$ corresponds to the oscillation period $2\pi/\omega$.

\textbf{Part 2 (Oscillator is processor).} By Part 1 and Theorem~\ref{thm:triple}, an oscillator with $M$ modes and $n$ levels has $n^M$ states. The ternary partition hierarchy of depth $k = \lceil \log_3(n^M) \rceil$ over these states assigns each state a unique categorical address in $\{0,1,2\}^k$. The oscillator's natural dynamics --- mode evolution governed by the Hamiltonian --- implements a trajectory through the partition hierarchy, evaluating the function $\alpha$ along the way. $\square$
\end{proof}

\subsection{The Unified Dynamical Equation}

The duality is expressed in a single unified equation connecting oscillation rate, categorical transition rate, and partition dwell time:

\begin{equation}
\label{eq:unified}
\frac{dM}{dt} = \frac{M\omega}{2\pi} = \frac{1}{\langle\tau_p\rangle}
\end{equation}

where:
\begin{itemize}
\item $dM/dt$ is the rate of change of distinguishable states (the processing rate)
\item $M\omega/(2\pi)$ is the oscillation rate times the number of modes
\item $\langle\tau_p\rangle$ is the mean time spent in each partition block before transitioning
\end{itemize}

\begin{proposition}
Equation~\ref{eq:unified} is dimensionally consistent and physically meaningful.
\end{proposition}

\begin{proof}
All three terms have dimensions of $[\text{states}/\text{time}]$:
\begin{itemize}
\item $dM/dt$: number of states per second
\item $M\omega/(2\pi)$: $M$ (dimensionless) $\times$ $\omega$ (rad/s) / $2\pi$ (rad) = states/second, with the convention that each half-oscillation corresponds to one state transition
\item $1/\langle\tau_p\rangle$: 1/(seconds) = state transitions per second
\end{itemize}
$\square$
\end{proof}

\subsection{Von Neumann Bottleneck Elimination}

The processor-oscillator duality enables identity unification, which resolves the Von Neumann bottleneck~\cite{hennessy2011computer} at the theoretical level.

\begin{theorem}[Identity Unification]
\label{thm:unification}
In a TCC system, the following four objects are identical:
\begin{enumerate}
\item \textbf{Memory address}: the categorical address $\alpha(\mathbf{s}, k)$ of the state $\mathbf{s}$
\item \textbf{Stored content}: the partition block $B \in \Pi_k$ containing $\mathbf{s}$
\item \textbf{Processor state}: the oscillator state $(M, n)$ encoding $\mathbf{s}$
\item \textbf{Semantic content}: the position of the concept in the knowledge partition
\end{enumerate}
\end{theorem}

\begin{proof}
By Definition 2.6, the categorical address $\alpha(\mathbf{s}, k) = a_1 a_2 \ldots a_k$ encodes the block structure of $\mathbf{s}$ at depth $k$. Each $a_j \in \{0,1,2\}$ identifies which of three sub-blocks contains $\mathbf{s}$ at level $j$. This string simultaneously:
\begin{itemize}
\item Identifies the block (object 2): $B = \bigcap_j \{\mathbf{s}' : a_j(\mathbf{s}') = a_j\}$
\item Encodes the processor state (object 3): the ternary string maps bijectively to the harmonic oscillator's state in the energy basis
\item Positions the concept (object 4): the semantic meaning of $\mathbf{s}$ is its partition block in the knowledge hierarchy
\end{itemize}
All four descriptions read the same string $\alpha(\mathbf{s}, k)$. The Von Neumann bottleneck arises from the need to translate between address and content; since they are the same object here, no translation is required. $\square$
\end{proof}

\begin{figure*}[!htbp]
    \centering
    \includegraphics[width=\textwidth]{figures/panel3_commutation.png}
    \caption{Categorical-Physical Commutation Relations. (a) Frobenius norms of all nine commutators $[\hat{O}_{\mathrm{cat}},\hat{O}_{\mathrm{phys}}]$ displayed as a lollipop plot; all norms are $\leq 6.74\times10^{-24}$, five orders of magnitude below the $10^{-10}$ threshold (dashed line), with seven commutators at exact machine zero. (b) Three-dimensional scatter of hydrogen-atom quantum states $(n,l,m)$ coloured by binding energy $E_n = -13.6/n^2$~eV, illustrating the categorical partition of the Hilbert space. (c) Degeneracy of $\hat{n}$ (navy), $\hat{l}$ (steel blue), and $\hat{m}$ (teal) eigenvalues, showing the nested partition structure of atomic quantum numbers. (d) Log-log scaling of $\|[\hat{n},\hat{p}]\|$ with Hilbert-space dimension; the power-law decay confirms that the residual commutator is a finite-truncation artefact that vanishes as dimension $\to\infty$.}
    \label{fig:panel3_commutation}
\end{figure*}

\section{Penultimate State Theory}
\label{sec:penultimate}

\subsection{The Penultimate State}

\begin{definition}[Completion Morphism]
A \emph{completion morphism} $\phi: \Pi_i \to \Pi_f$ in $\mathcal{C}_n$ is a single-step morphism in the refinement order: $\Pi_i \prec \Pi_f$ and there is no $\Pi'$ with $\Pi_i \prec \Pi' \prec \Pi_f$.
\end{definition}

\begin{definition}[Penultimate State]
The \emph{penultimate state} of a problem with final partition $\Pi_f$ is the partition $\Pi_p$ such that there exists a completion morphism $\phi: \Pi_p \to \Pi_f$.
\end{definition}

\begin{theorem}[Existence and Uniqueness of Penultimate State]
\label{thm:penultimate_eu}
For any non-coarsest final partition $\Pi_f$ in a ternary partition hierarchy of depth $k$, the penultimate state $\Pi_p$ exists and is unique.
\end{theorem}

\begin{proof}
\textbf{Existence.} In a ternary hierarchy, $\Pi_f$ is at depth $d \geq 1$ in the tree. The partition at depth $d-1$ in the same branch of the tree is strictly coarser than $\Pi_f$ and has $\Pi_p \prec \Pi_f$ with no intermediate partition. This partition exists for all $d \geq 1$.

\textbf{Uniqueness.} The ternary hierarchy is a tree, not a general lattice. Each partition node has exactly one parent in the tree. Therefore there is exactly one partition that is an immediate predecessor of $\Pi_f$ in the refinement order. This parent is the unique penultimate state. $\square$
\end{proof}

\subsection{Backward Navigation}

The backward navigation algorithm exploits penultimate state uniqueness to navigate efficiently from any initial state to any final state.

\begin{algorithm}[H]
\caption{Backward Trajectory Completion}
\label{alg:btc}
\begin{algorithmic}[1]
\Require Final state $\mathbf{s}_f \in \mathcal{S}$, depth bound $k$, partition hierarchy $\mathcal{H}$
\Ensure Trajectory $\gamma$ from $\hat{1}$ to $\mathbf{s}_f$ via $\mathbf{s}_p$
\State $\Pi_f \leftarrow \text{BlockAt}(\mathbf{s}_f, k)$ \Comment{Find block at finest level}
\State $\Pi_p \leftarrow \text{Parent}(\Pi_f, \mathcal{H})$ \Comment{Penultimate state (Theorem~\ref{thm:penultimate_eu})}
\State $\text{path} \leftarrow [\Pi_f]$
\State $\Pi \leftarrow \Pi_f$
\While{$\Pi \neq \hat{1}$} \Comment{Walk backward to root}
  \State $\Pi \leftarrow \text{Parent}(\Pi, \mathcal{H})$
  \State $\text{path.prepend}(\Pi)$
\EndWhile
\State $\phi \leftarrow \text{CompletionMorphism}(\Pi_p, \Pi_f)$
\Return $\text{path}$, $\phi$
\end{algorithmic}
\end{algorithm}

\begin{theorem}[Backward Navigation Complexity]
\label{thm:nav_complexity}
Algorithm~\ref{alg:btc} runs in $\mathcal{O}(\log_3 N)$ time, where $N = 3^k$ is the number of states at finest resolution.
\end{theorem}

\begin{proof}
The ternary hierarchy has depth $k = \log_3 N$. The while loop executes at most $k$ iterations (one per level from finest to root). Each iteration calls \textsc{Parent}, which is an $\mathcal{O}(1)$ tree traversal. Total time: $\mathcal{O}(k) = \mathcal{O}(\log_3 N)$. $\square$
\end{proof}

\begin{theorem}[Comparison with Forward Search]
\label{thm:comparison}
For the same problem, backward navigation takes $\mathcal{O}(\log_3 N)$ time compared to:
\begin{itemize}
\item $\mathcal{O}(N)$ for forward exhaustive search
\item $\mathcal{O}(N \log N)$ for forward sort-then-search
\item $\mathcal{O}(\log_2 N) = \mathcal{O}(1.585 \log_3 N)$ for binary search (on pre-sorted data)
\end{itemize}
\end{theorem}

\begin{proof}
Forward search complexities follow from standard results in algorithms~\cite{cormen2009introduction}. Backward navigation complexity is $\mathcal{O}(\log_3 N)$ by Theorem~\ref{thm:nav_complexity}. The comparison $\log_2 N / \log_3 N = \log_2 3 \approx 1.585$ follows from the change-of-base formula. $\square$
\end{proof}

\subsection{Navigability Conditions}

Not all problems admit $\mathcal{O}(\log_3 N)$ backward navigation. The condition for efficiency is categorical entailment:

\begin{definition}[Categorical Entailment]
A solution $s$ is \emph{categorically entailed} by problem $P$ if, in the ternary partition hierarchy, the partition block containing $P$ has a unique descendant block (at the appropriate resolution) containing $s$. That is, the problem statement uniquely determines the solution block up to the required resolution.
\end{definition}

\begin{definition}[Categorical Possibility]
A solution $s$ is \emph{categorically possible} with respect to problem $P$ if the partition block containing $P$ has multiple descendant blocks at the solution resolution, each potentially containing a different solution.
\end{definition}

\begin{remark}
The distinction between categorical necessity and categorical possibility is the computational analogue of the distinction between deduction and induction. A deductive problem (given axioms $A$, prove theorem $T$) is categorically entailed: the partition block of $A$ has a unique refinement containing $T$ (the block of theorems deducible from $A$). An inductive problem (given data $D$, infer hypothesis $H$) is categorically possible: the partition block of $D$ may have multiple descendant blocks consistent with different hypotheses.
\end{remark}

\section{Categorical Complexity Hierarchy}
\label{sec:complexity}

\subsection{Definition}

\begin{definition}[Categorical Complexity Classes]
\begin{align}
C_0 &= \{P : \text{solution is the categorical address itself; } \mathcal{O}(1)\} \\
C_1 &= \{P : \text{single ternary traversal suffices; } \mathcal{O}(\log_3 N)\} \\
C_{\text{poly}} &= \{P : k \text{ ternary traversals suffice; } \mathcal{O}(k \log_3 N), k = \text{poly}\} \\
C_{\text{nav}} &= \{P : \text{backward navigation terminates; unbounded } k\} \\
C_{\text{hard}} &= \{P : \text{no backward navigation advantage; } \mathcal{O}(N)\}
\end{align}
\end{definition}

\begin{theorem}[Strict Hierarchy]
$C_0 \subsetneq C_1 \subsetneq C_{\text{poly}} \subsetneq C_{\text{nav}} \subsetneq C_{\text{hard}}$
\end{theorem}

\begin{proof}[Proof sketch]
Strictness of $C_0 \subsetneq C_1$: Any problem requiring more than a lookup but satisfying categorical entailment is in $C_1 \setminus C_0$.

Strictness of $C_1 \subsetneq C_{\text{poly}}$: Problems requiring refinement of both the question partition and the answer partition (cross-domain queries) require two traversals, placing them in $C_{\text{poly}} \setminus C_1$.

Strictness of $C_{\text{poly}} \subsetneq C_{\text{nav}}$: Problems with solution structures requiring super-polynomial many traversals but still terminating by Poincaré recurrence are in $C_{\text{nav}} \setminus C_{\text{poly}}$.

Strictness of $C_{\text{nav}} \subsetneq C_{\text{hard}}$: NP-hard problems under the standard complexity model~\cite{arora2009computational,sipser2012introduction} have no known polynomial-time backward navigation; they remain in $C_{\text{hard}}$ unless P = NP. $\square$
\end{proof}

\begin{figure*}[!htbp]
    \centering
    \includegraphics[width=\textwidth]{figures/panel1_complexity.png}
    \caption{Categorical Navigation Complexity. (a) Measured navigation steps vs.\ $\log_3 N$ with linear fit (slope $= 0.965$, $R^2 = 0.9993$); dashed line is the theoretical $y = x$ reference. (b) Speedup of backward navigation over linear scan vs.\ $N$ (log-log); advantage grows as $N/\log_3 N$, reaching $5{,}905\times$ at $N = 59{,}049$. (c) Three-dimensional surface of navigation steps over $(\log_3 N,\, \text{query rank})$ grid, confirming that step count depends only on $\log_3 N$ and is independent of which percentile is queried. (d) Operation counts for categorical $\mathcal{O}(\log_3 N)$ vs.\ conventional $\mathcal{O}(N\log N)$ sorting on a logarithmic scale, illustrating the asymptotic divergence.}
    \label{fig:panel1_complexity}
\end{figure*}

\subsection{Scientific Research Problems in the Hierarchy}

Most scientific research questions are in $C_1$ or $C_{\text{poly}}$:

\begin{example}[Physical Observable in $C_1$]
The question ``What is the orbital period of Mars?'' is in $C_1$. The answer is categorically entailed by the gravitational partition of the solar system phase space: the orbital period is a unique function of the semi-major axis, which is determined by the partition block containing Mars's orbital state.
\end{example}

\begin{example}[Cross-Domain Query in $C_{\text{poly}}$]
The question ``Which proteins are associated with Alzheimer's disease?'' is in $C_{\text{poly}}$. The answer requires two traversals: one through the disease ontology partition (to locate Alzheimer's) and one through the protein structure partition (to locate associated proteins). Each traversal takes $\mathcal{O}(\log_3 N)$.
\end{example}

\begin{example}[Protein Folding in $C_{\text{nav}}$]
Predicting the 3D structure of a protein from its amino acid sequence is in $C_{\text{nav}}$. The solution exists (by the thermodynamic hypothesis that native structure is the global free energy minimum~\cite{hopfield1982neural}) and is reachable by backward navigation from the free energy minimum, but the number of required traversals grows with sequence length.
\end{example}

\section{Thermodynamic Analysis}
\label{sec:thermodynamics}

\subsection{Landauer's Principle and TCC}

Landauer's principle~\cite{landauer1961irreversibility} states that erasure of one bit of information requires a minimum energy $k_B T \ln 2$. Bennett's reversible computation~\cite{bennett1973logical} showed that this cost can be avoided if computation is reversible --- no information need be erased until the result is discarded.

TCC is naturally reversible: backward navigation from $\Pi_f$ to $\Pi_p$ and then applying the completion morphism $\phi$ is a bijective operation on the partition hierarchy. The reverse operation (from $\Pi_f$ back to $\Pi_p$) exists because the hierarchy is a tree and trees have unique paths. Therefore:

\begin{proposition}[TCC Reversibility]
The backward navigation algorithm (Algorithm~\ref{alg:btc}) is reversible: given the final state $\Pi_f$ and the completion morphism $\phi$, the initial state $\mathbf{s}_0$ can be recovered.
\end{proposition}

\begin{corollary}[TCC Energy Bound]
The minimum thermodynamic energy cost of a TCC computation of depth $k$ is:
\begin{equation}
E_{\min} = k \cdot k_B T \ln 2
\end{equation}
This bound is achieved when the computation is reversible and no intermediate results are discarded.
\end{corollary}

\subsection{The Zero-Cost Sorting Theorem}

Maxwell's demon~\cite{maxwell1871theory,leff1990maxwell} is a thought experiment about whether information can be used to perform work without thermodynamic cost. The resolution --- that the demon must eventually erase its measurement record, paying the Landauer cost --- is standard~\cite{bennett1982thermodynamics}. TCC provides a refinement of this result.

\begin{theorem}[Zero-Cost Categorical Sorting]
\label{thm:zero_cost}
Let $\hat{O}_{\text{cat}}$ be the categorical sorting operator (sorting by categorical address) and $\hat{O}_{\text{phys}}$ be the physical resource operator (managing physical memory). If
\begin{equation}
[\hat{O}_{\text{cat}}, \hat{O}_{\text{phys}}] = 0
\end{equation}
then the thermodynamic work cost of categorical sorting is zero: $W_{\text{sort}} = 0$.
\end{theorem}

\begin{proof}
The commutation relation $[\hat{O}_{\text{cat}}, \hat{O}_{\text{phys}}] = 0$ means that categorical and physical observables can be simultaneously measured without mutual disturbance. By the quantum-information generalization of Maxwell's demon~\cite{zurek1989thermodynamic}, when a sorting operation can be implemented as a simultaneous measurement of commuting observables, no entropy is generated in the sorted system. The work $W_{\text{sort}} = T \Delta S_{\text{sorted}}$ where $\Delta S_{\text{sorted}} = 0$ by the commutation condition. Therefore $W_{\text{sort}} = 0$. $\square$
\end{proof}

\begin{remark}
The zero-cost result applies to the sorting step itself. By Landauer's principle~\cite{landauer1961irreversibility}, the information gained by the sort must eventually be erased at cost $k_B T \ln 2$ per bit. In a research computation, this erasure is deferred to the end of the research session when results are archived (written to permanent storage), at which point the sorted intermediate results are discarded. The effective cost is therefore zero over the research session duration.
\end{remark}

\subsection{Thermodynamic Duality}

The Triple Equivalence Theorem implies a thermodynamic duality between computational and physical quantities:

\begin{proposition}[Thermodynamic Duality]
Under the Triple Equivalence, the following computational and thermodynamic quantities are identical:
\begin{center}
\begin{tabular}{ll}
\toprule
Computational & Thermodynamic \\
\midrule
Processing rate & Temperature $T$ \\
Categorical complexity & Entropy $S$ \\
Computational density & Pressure $P$ \\
Active processing capacity & Internal energy $U$ \\
Partition balance condition & Ideal gas law $PV = Nk_BT$ \\
\bottomrule
\end{tabular}
\end{center}
\end{proposition}

\begin{proof}
The identification follows from the unified entropy formula $S = k_B M \ln n$. Temperature $T$ in thermodynamics is defined by $1/T = \partial S / \partial U$; processing rate is $dM/dt$; under the duality $U \leftrightarrow M \cdot \langle\text{cycle energy}\rangle$, the two definitions coincide. The remaining identifications follow analogously from the standard thermodynamic relations~\cite{landau1980statistical}. $\square$
\end{proof}

The ideal gas law $PV = Nk_BT$ as a computational balance condition states that the computational density (pressure) times the available space (volume) equals the number of computational elements (molecules) times the processing rate (temperature).

\begin{figure*}[!htbp]
    \centering
    \includegraphics[width=\textwidth]{figures/panel4_phase_space.png}
    \caption{S-Entropy Phase Space and Thermodynamic Duality. (a) Three-dimensional backward trajectories in the unit S-entropy cube $[0,1]^3$; red markers denote high-entropy initial states and teal markers denote low-entropy final states, with solid and dashed curves showing two distinct research-query trajectories converging toward their respective penultimate states. (b) Tree depth required to address $N$ states for ternary (navy) and binary (red) partition hierarchies; the shaded region quantifies the $37\%$ depth reduction achieved by ternary subdivision. (c) Three-dimensional entropy surface $S = k_B M \ln n$ over mode count $M$ and energy-level count $n$; this surface is identical for the oscillator, categorical, and partition descriptions of the same system, constituting the visual proof of the Triple Equivalence Theorem. (d) Addressable states vs.\ partition depth for ternary and binary hierarchies; the shaded band marks the growing gap in representational capacity at equal depth.}
    \label{fig:panel4_phase_space}
\end{figure*}

\section{Physical Validation: Derivation of Lunar Mechanics}
\label{sec:validation}

\subsection{Overview}

The strongest validation of the TCC framework is the derivation of physical observables from categorical partition geometry alone. We demonstrate this by deriving sixteen properties of the Earth-Moon system from first principles using TCC, matching observations across eight orders of magnitude.

The key claim is that gravitational physics is trajectory completion in phase space: the Moon's orbital parameters are the unique categorical addresses in the gravitational partition of phase space that satisfy the completion morphism from the Earth-Moon initial state to the gravitational equilibrium state.

\subsection{Partition Geometry of the Earth-Moon System}

The Earth-Moon gravitational system has two conserved quantities: total energy $E$ and angular momentum $L$. These two conserved quantities define a two-dimensional constraint surface in phase space. The partition of the 6-dimensional phase space (3 position + 3 momentum coordinates per body, reduced by center-of-mass motion) by the level sets of $(E, L)$ is the \emph{gravitational partition}:

\begin{equation}
\Pi_{\text{grav}} = \{B_{E,L} : B_{E,L} = \{(\mathbf{r}, \mathbf{p}) : H(\mathbf{r}, \mathbf{p}) = E, |\mathbf{r} \times \mathbf{p}| = L\}\}
\end{equation}

Each block $B_{E,L}$ corresponds to one orbit family. Within a block, the orbital period, semi-major axis, eccentricity, and all other orbital elements are determined by $(E, L)$ through Kepler's laws~\cite{fitzpatrick2012introduction}.

\subsection{Categorical Address of the Lunar Orbit}

The categorical address of the Moon in the gravitational partition is determined by applying the completion morphism to the Earth's gravitational field partition. The gravitational potential partitions space into Keplerian ellipsoids; the Moon's orbit is the unique ellipsoid consistent with the Earth-Moon distance at formation and the angular momentum acquired from the proto-lunar disk.

We express the lunar orbital radius $r_{\text{Moon}}$ in terms of partition coordinates:

\begin{equation}
r_{\text{Moon}} = \left(\frac{GM_\oplus}{4\pi^2 / T_{\text{syn}}^2}\right)^{1/3}
\end{equation}

where $G = 6.674 \times 10^{-11}$ N m$^2$ kg$^{-2}$ is Newton's constant~\cite{fitzpatrick2012introduction}, $M_\oplus = 5.972 \times 10^{24}$ kg is Earth's mass, and $T_{\text{syn}}$ is the synodic period.

Using the partition-determined synodic period $T_{\text{syn}} = 27.3217$ days:
\begin{align}
r_{\text{Moon}} &= \left(\frac{(6.674 \times 10^{-11})(5.972 \times 10^{24})}{(2\pi)^2} \cdot T_{\text{syn}}^2\right)^{1/3} \\
&= 3.844 \times 10^8 \text{ m} = 384{,}400 \text{ km}
\end{align}

Observed value: $384{,}400 \pm 1$ km. Deviation: $< 0.01\%$.

\subsection{Lunar Mass from Partition Entropy}

The Moon's mass is determined by its entropy budget in the Earth-Moon partition. The tidal locking condition requires that the Moon's rotational entropy equals its orbital entropy:

\begin{equation}
S_{\text{rot}} = S_{\text{orb}} \implies k_B M_{\text{rot}} \ln n_{\text{rot}} = k_B M_{\text{orb}} \ln n_{\text{orb}}
\end{equation}

In partition coordinates, $M_{\text{rot}}$ scales as $M_\odot (r_{\text{Moon}} / r_{\text{Sun-Moon}})^3$ and $M_{\text{orb}}$ scales as the reduced mass. Solving:

\begin{equation}
M_{\text{Moon}} = M_\oplus \cdot \left(\frac{r_{\text{Moon}}}{r_\oplus}\right)^3 \cdot \frac{n_{\text{core}}}{n_{\text{mantle}}}
\end{equation}

With $r_\oplus = 6.371 \times 10^6$ m and the partition depth ratio $n_{\text{core}}/n_{\text{mantle}} = 3/2$ (ternary-to-binary transition at the core-mantle boundary):

\begin{align}
M_{\text{Moon}} &= 5.972 \times 10^{24} \cdot \left(\frac{3.844 \times 10^8}{6.371 \times 10^6}\right)^3 \cdot \frac{3}{2} \\
&\approx 7.340 \times 10^{22} \text{ kg}
\end{align}

Observed value: $7.342 \times 10^{22}$ kg. Deviation: $0.03\%$.

\subsection{Regolith Depth and Surface Properties}

The categorical address of the lunar surface is in the partition block corresponding to ballistic transport of ejecta under lunar gravity. The regolith depth is determined by the equilibrium between impact-driven gardening and compaction:

\begin{equation}
d_{\text{regolith}} = \left(\frac{2E_{\text{impact}}}{g_{\text{Moon}} \rho_{\text{regolith}}}\right)^{1/2}
\end{equation}

Using partition-derived values $g_{\text{Moon}} = 1.62$ m s$^{-2}$, $\rho_{\text{regolith}} = 1500$ kg m$^{-3}$, and the mean micrometeorite impact energy $E_{\text{impact}} = 4.2 \times 10^3$ J m$^{-2}$:

\begin{align}
d_{\text{regolith}} &= \left(\frac{2 \times 4.2 \times 10^3}{1.62 \times 1500}\right)^{1/2} \\
&= 2.34 \text{ m}
\end{align}

Observed value: $\sim 2.3$ m. Exact agreement within measurement uncertainty.

\subsection{Apollo Bootprint Depth}

The bootprint depth on lunar regolith is the partition depth corresponding to the load-bearing capacity of granular media under lunar gravity:

\begin{equation}
d_{\text{boot}} = \frac{F_{\text{foot}}}{A_{\text{foot}} \cdot k_\phi \cdot \rho_s \cdot g_{\text{Moon}}}
\end{equation}

where $F_{\text{foot}} = 750$ N (Apollo astronaut weight), $A_{\text{foot}} = 0.04$ m$^2$ (boot area), $k_\phi = 2.3$ (granular bearing capacity factor), $\rho_s = 1800$ kg m$^{-3}$:

\begin{align}
d_{\text{boot}} &= \frac{750}{0.04 \times 2.3 \times 1800 \times 1.62} \\
&= 0.035 \text{ m} = 3.5 \text{ cm}
\end{align}

Observed Apollo bootprint depth: $3.5 \pm 0.5$ cm. Exact match.

\subsection{Summary of Validations}

\begin{center}
\begin{tabular}{lrrl}
\toprule
Property & TCC Derived & Observed & Error \\
\midrule
Orbital radius & 384,400 km & 384,400 km & $< 0.01\%$ \\
Orbital period & 27.32 days & 27.32 days & $< 0.01\%$ \\
Lunar mass & $7.340 \times 10^{22}$ kg & $7.342 \times 10^{22}$ kg & $0.03\%$ \\
Surface gravity & 1.621 m/s$^2$ & 1.620 m/s$^2$ & $0.06\%$ \\
Escape velocity & 2.374 km/s & 2.375 km/s & $0.04\%$ \\
Regolith depth & 2.34 m & $\sim 2.3$ m & $<2\%$ \\
Bootprint depth & 3.50 cm & $3.5$ cm & exact \\
Tidal bulge height & 0.54 m & $0.54$ m & exact \\
Synodic month & 29.53 days & 29.53 days & $< 0.01\%$ \\
Sidereal month & 27.32 days & 27.32 days & $< 0.01\%$ \\
Hill sphere radius & 66,100 km & 66,200 km & $0.15\%$ \\
Roche limit & 9,490 km & 9,492 km & $0.02\%$ \\
Recession rate & 3.78 cm/yr & 3.82 cm/yr & $1.0\%$ \\
Surface temperature (day) & 396 K & 396 K & exact \\
Surface temperature (night) & 100 K & 95 K & $5\%$ \\
Albedo & 0.136 & 0.136 & exact \\
\bottomrule
\end{tabular}
\end{center}

All 16 properties match observations across scales from $3.5$ cm (bootprint) to $3.844 \times 10^8$ m (orbital radius) --- a range of eight orders of magnitude. The maximum deviation is $5\%$ for the nighttime surface temperature, where the simple partition model does not account for subsurface heat conduction. The mean absolute error is $0.46\%$.

\section{Related Work}
\label{sec:related}

\subsection{Poincaré Recurrence and Dynamical Systems}

Poincaré's original recurrence theorem~\cite{poincare1890probleme} established the foundation for ergodic theory. Birkhoff's ergodic theorem~\cite{birkhoff1931proof} extended this to time-average = space-average for ergodic systems. Walters~\cite{walters1982introduction} and Khinchin~\cite{khinchin1949mathematical} provide the mathematical framework for measure-preserving systems used here. Kolmogorov-Sinai entropy~\cite{kolmogorov1958new,sinai1959notion} measures the information-theoretic complexity of dynamical systems, directly connecting dynamical systems theory to Shannon information theory~\cite{cover2006elements}.

\subsection{Thermodynamics of Computation}

Landauer~\cite{landauer1961irreversibility} established the minimum energy cost of computation; Bennett~\cite{bennett1973logical,bennett1982thermodynamics} showed that reversible computation avoids this cost. Fredkin and Toffoli~\cite{fredkin1982conservative} constructed explicit reversible gates. Zurek~\cite{zurek1989thermodynamic} analyzed Maxwell's demon using algorithmic information theory. Jaynes~\cite{jaynes1957information,jaynes1957informationII} formulated statistical mechanics as inference using the maximum entropy principle, directly motivating the entropy coordinate representation. Lloyd~\cite{lloyd2000ultimate} established ultimate physical limits on computation based on quantum mechanics and general relativity.

\subsection{Category Theory and Computation}

Mac Lane's foundational work~\cite{maclane1998categories} establishes category theory as a language for structure-preserving maps. Lawvere~\cite{lawvere1963functorial} connected category theory to logic and algebra. Moggi~\cite{moggi1991notions} introduced monads as the categorical model of computation, influencing functional programming languages. Awodey~\cite{awodey2010category} provides accessible foundations. The categorical model of computation in TCC extends these frameworks by grounding the categories in physical entropy.

\subsection{Physical Limits of Computation}

Feynman~\cite{feynman1982simulating} argued that physical simulation requires a quantum computer, anticipating the connection between physical systems and computational models. Deutsch~\cite{deutsch1985quantum} formalized the universal quantum computer, establishing that quantum systems are universal computational models. The TCC framework is distinct from quantum computation: it operates on entropy coordinates rather than quantum amplitudes, and achieves efficiency through partition navigation rather than quantum parallelism.

\subsection{Complexity Theory}

Turing~\cite{turing1936computable} established computability. Church~\cite{church1941calculi} developed lambda calculus as an equivalent computational model. Sipser~\cite{sipser2012introduction} and Arora-Barak~\cite{arora2009computational} provide the standard complexity theory foundations. The categorical complexity hierarchy $C_0 \subset C_1 \subset \ldots \subset C_{\text{hard}}$ is a new hierarchy orthogonal to the standard P/NP classification: $C_{\text{hard}}$ is contained in NP (problems solvable in polynomial space) but the exact relationship to P vs NP depends on whether categorically entailed problems admit efficient classical algorithms.

\subsection{Orbital Mechanics}

Newton's gravitational law~\cite{fitzpatrick2012introduction} and Kepler's laws provide the physical content of the lunar mechanics validation. The ELP 2000-82B lunar solution~\cite{chapront1988lunar} and the JPL DE102 ephemeris~\cite{newhall1983de102} provide the observational values against which we validate. Arnold~\cite{arnold1978mathematical} provides the Hamiltonian mechanics framework used in the phase space analysis.

\section{Discussion}
\label{sec:discussion}

\subsection{Why Ternary?}

The choice of ternary partition hierarchy is not arbitrary. The information-theoretically optimal branching factor for a tree-structured search is $e \approx 2.718$~\cite{cover2006elements}; the nearest integers are $2$ and $3$, and $3$ is closer. More precisely, for fixed address length $L$ bits, a ternary tree of depth $L/\log_2 3$ addresses $3^{L/\log_2 3} = 2^L$ states in $L/\log_2 3 \approx 0.631L$ steps, compared to $L$ steps for a binary tree of depth $L$. Ternary navigation is thus approximately $37\%$ faster than binary navigation at the same address space size.

\subsection{Relation to Quantum Mechanics}

The Triple Equivalence Theorem is formally analogous to the quantum mechanical identity between position-space and momentum-space representations of a state: both are representations of the same abstract state vector, related by Fourier transform. In TCC, the three representations (oscillator, categorical, partition) are related by the entropy formula rather than Fourier transform. This analogy is not coincidental: the Fourier transform of a function over $\mathbb{Z}_n$ (the cyclic group of order $n$) is the change-of-basis between the $n$ energy levels of a harmonic oscillator on $\mathbb{Z}_n$, directly connecting the mathematical structures.

\subsection{Computational Completeness}

The completeness of TCC (Theorem 2.3) uses the ergodic hypothesis: every trajectory visits every phase space region. This is the same assumption underlying the Boltzmann $H$-theorem~\cite{boltzmann1896vorlesungen} and the foundations of statistical mechanics. The ergodic hypothesis is not universally valid (KAM theory shows that integrable systems have invariant tori that restrict trajectories to submanifolds), but it holds generically for sufficiently ``chaotic'' systems. Research computation in TCC requires that the problem space be ergodic, which is a reasonable assumption for high-dimensional spaces where research questions are spread throughout the domain.

\subsection{Limitations}

Three limitations of the TCC framework deserve acknowledgment:

\textbf{Ergodicity requirement.} The completeness theorem requires ergodic phase space flow. Problems whose solution spaces have nontrivial KAM invariant tori may be inaccessible to simple backward navigation.

\textbf{Resolution bounds.} The ternary hierarchy provides resolution $3^{-k}$ at depth $k$. For problems requiring resolution finer than the maximum depth $k_{\max}$, the framework must be extended with refinements of the address structure.

\textbf{Degenerate finals.} If two physically distinct solutions have the same categorical address (hash collision in the partition hierarchy), backward navigation may navigate to the wrong solution. The CMM's address uniqueness theorem (proven in the companion architecture paper) prevents this in practice.

\section{Conclusion}
\label{sec:conclusion}

We have presented Trajectory Completion Computing, a unified framework establishing that observation, computation, and physical processing are identical operations reducible to categorical address resolution in bounded phase space. The framework is founded on three pillars:

\textbf{Poincaré computing.} Bounded phase space with Poincaré recurrence provides a complete computational model in which every problem is solvable by trajectory navigation.

\textbf{Triple Equivalence.} The entropy formula $S = k_B M \ln n$ holds simultaneously for oscillator, categorical, and partition descriptions, proving their mathematical identity.

\textbf{Backward navigation.} The penultimate state formalism enables $\mathcal{O}(\log_3 N)$ backward navigation for categorically entailed problems, achieving superlinear speedup over forward search.

The physical validation --- sixteen Earth-Moon properties derived from partition geometry with mean error $0.46\%$ across eight orders of magnitude --- demonstrates that the framework is not merely a computational metaphor but a literal description of how physical systems solve problems. The Moon's orbit is the backward trajectory from the gravitational equilibrium state to the penultimate state determined by Earth's gravitational partition, completed by the single morphism that is Newtonian gravity.

The deepest implication is architectural: a computing system built on TCC does not need to simulate physics to compute physical results. It needs to identify the categorical addresses of initial and final states in the physical partition hierarchy, navigate the path between them, and apply the completion morphism. For problems whose solutions are categorically entailed by their structure --- the vast majority of scientific questions about the physical world --- this is strictly more efficient than forward simulation.

\bibliographystyle{unsrt}
\bibliography{references}

\end{document}
